% Perbandingan jabat tangan TLS 1.2 dan TLS 1.3 (diagram urutan).
% Dirujuk dari section-04.tex dengan % Perbandingan jabat tangan TLS 1.2 dan TLS 1.3 (diagram urutan).
% Dirujuk dari section-04.tex dengan % Perbandingan jabat tangan TLS 1.2 dan TLS 1.3 (diagram urutan).
% Dirujuk dari section-04.tex dengan % Perbandingan jabat tangan TLS 1.2 dan TLS 1.3 (diagram urutan).
% Dirujuk dari section-04.tex dengan \input{figures/jabat-tangan-tls}.
\begin{figure}[htbp]
    \centering
    \begin{tikzpicture}[
        panah/.style={-{Latex[length=2.2mm]}, thick},
        balik/.style={-{Latex[length=2.2mm]}, thick, dashed},
        label/.style={font=\scriptsize, align=center, fill=white, inner sep=1.5pt},
        kepala/.style={draw, rounded corners, fill=blue!6, font=\small,
                       minimum width=1.7cm, minimum height=0.6cm},
        hidup/.style={dashed, gray},
    ]

    % ================= TLS 1.2 (kiri) =================
    \node[font=\small\bfseries] at (2.8,0.95) {TLS 1.2: dua putaran};
    \node[kepala] (k12) at (0,0.35) {Klien};
    \node[kepala] (s12) at (5.6,0.35) {Server};
    \draw[hidup] (0,0) -- (0,-7.2);
    \draw[hidup] (5.6,0) -- (5.6,-7.2);

    \draw[panah] (0,-1.3) -- node[label] {ClientHello} (5.6,-1.3);
    \draw[panah] (5.6,-2.7) -- node[label]
        {ServerHello + Certificate\\ + ServerKeyExchange\\ + ServerHelloDone} (0,-2.7);
    \draw[panah] (0,-4.2) -- node[label]
        {ClientKeyExchange\\ + ChangeCipherSpec\\ + Finished} (5.6,-4.2);
    \draw[panah] (5.6,-5.4) -- node[label]
        {ChangeCipherSpec\\ + Finished} (0,-5.4);
    \draw[balik] (0,-6.6) -- node[label] {Data terenkripsi} (5.6,-6.6);

    % ================= TLS 1.3 (kanan) =================
    \begin{scope}[xshift=7.6cm]
    \node[font=\small\bfseries] at (2.8,0.95) {TLS 1.3: satu putaran};
    \node[kepala] (k13) at (0,0.35) {Klien};
    \node[kepala] (s13) at (5.6,0.35) {Server};
    \draw[hidup] (0,0) -- (0,-7.2);
    \draw[hidup] (5.6,0) -- (5.6,-7.2);

    \draw[panah] (0,-1.3) -- node[label] {ClientHello\\ + key share} (5.6,-1.3);
    \draw[panah] (5.6,-3.4) -- node[label]
        {ServerHello\\ + key share\\ \{Certificate,\\ CertificateVerify,\\ Finished\}} (0,-3.4);
    \draw[panah] (0,-5.4) -- node[label] {Finished} (5.6,-5.4);
    \draw[balik] (0,-6.6) -- node[label] {Data terenkripsi} (5.6,-6.6);
    \end{scope}

    \end{tikzpicture}
    \caption{Perbandingan jabat tangan TLS 1.2 dan TLS 1.3. TLS 1.2 menuntut
    dua putaran penuh sebelum data dapat dikirim, karena kunci sesi baru
    ditentukan setelah klien menerima sertifikat server. TLS 1.3 memangkasnya
    menjadi satu putaran: klien sudah mengirim bagian kunci publiknya pada
    pesan pertama, sehingga server dapat menghitung kunci sesi lebih awal dan
    seluruh pesan setelah \textit{ServerHello} sudah terenkripsi}
    \label{fig:jabat-tangan-tls}
\end{figure}
.
\begin{figure}[htbp]
    \centering
    \begin{tikzpicture}[
        panah/.style={-{Latex[length=2.2mm]}, thick},
        balik/.style={-{Latex[length=2.2mm]}, thick, dashed},
        label/.style={font=\scriptsize, align=center, fill=white, inner sep=1.5pt},
        kepala/.style={draw, rounded corners, fill=blue!6, font=\small,
                       minimum width=1.7cm, minimum height=0.6cm},
        hidup/.style={dashed, gray},
    ]

    % ================= TLS 1.2 (kiri) =================
    \node[font=\small\bfseries] at (2.8,0.95) {TLS 1.2: dua putaran};
    \node[kepala] (k12) at (0,0.35) {Klien};
    \node[kepala] (s12) at (5.6,0.35) {Server};
    \draw[hidup] (0,0) -- (0,-7.2);
    \draw[hidup] (5.6,0) -- (5.6,-7.2);

    \draw[panah] (0,-1.3) -- node[label] {ClientHello} (5.6,-1.3);
    \draw[panah] (5.6,-2.7) -- node[label]
        {ServerHello + Certificate\\ + ServerKeyExchange\\ + ServerHelloDone} (0,-2.7);
    \draw[panah] (0,-4.2) -- node[label]
        {ClientKeyExchange\\ + ChangeCipherSpec\\ + Finished} (5.6,-4.2);
    \draw[panah] (5.6,-5.4) -- node[label]
        {ChangeCipherSpec\\ + Finished} (0,-5.4);
    \draw[balik] (0,-6.6) -- node[label] {Data terenkripsi} (5.6,-6.6);

    % ================= TLS 1.3 (kanan) =================
    \begin{scope}[xshift=7.6cm]
    \node[font=\small\bfseries] at (2.8,0.95) {TLS 1.3: satu putaran};
    \node[kepala] (k13) at (0,0.35) {Klien};
    \node[kepala] (s13) at (5.6,0.35) {Server};
    \draw[hidup] (0,0) -- (0,-7.2);
    \draw[hidup] (5.6,0) -- (5.6,-7.2);

    \draw[panah] (0,-1.3) -- node[label] {ClientHello\\ + key share} (5.6,-1.3);
    \draw[panah] (5.6,-3.4) -- node[label]
        {ServerHello\\ + key share\\ \{Certificate,\\ CertificateVerify,\\ Finished\}} (0,-3.4);
    \draw[panah] (0,-5.4) -- node[label] {Finished} (5.6,-5.4);
    \draw[balik] (0,-6.6) -- node[label] {Data terenkripsi} (5.6,-6.6);
    \end{scope}

    \end{tikzpicture}
    \caption{Perbandingan jabat tangan TLS 1.2 dan TLS 1.3. TLS 1.2 menuntut
    dua putaran penuh sebelum data dapat dikirim, karena kunci sesi baru
    ditentukan setelah klien menerima sertifikat server. TLS 1.3 memangkasnya
    menjadi satu putaran: klien sudah mengirim bagian kunci publiknya pada
    pesan pertama, sehingga server dapat menghitung kunci sesi lebih awal dan
    seluruh pesan setelah \textit{ServerHello} sudah terenkripsi}
    \label{fig:jabat-tangan-tls}
\end{figure}
.
\begin{figure}[htbp]
    \centering
    \begin{tikzpicture}[
        panah/.style={-{Latex[length=2.2mm]}, thick},
        balik/.style={-{Latex[length=2.2mm]}, thick, dashed},
        label/.style={font=\scriptsize, align=center, fill=white, inner sep=1.5pt},
        kepala/.style={draw, rounded corners, fill=blue!6, font=\small,
                       minimum width=1.7cm, minimum height=0.6cm},
        hidup/.style={dashed, gray},
    ]

    % ================= TLS 1.2 (kiri) =================
    \node[font=\small\bfseries] at (2.8,0.95) {TLS 1.2: dua putaran};
    \node[kepala] (k12) at (0,0.35) {Klien};
    \node[kepala] (s12) at (5.6,0.35) {Server};
    \draw[hidup] (0,0) -- (0,-7.2);
    \draw[hidup] (5.6,0) -- (5.6,-7.2);

    \draw[panah] (0,-1.3) -- node[label] {ClientHello} (5.6,-1.3);
    \draw[panah] (5.6,-2.7) -- node[label]
        {ServerHello + Certificate\\ + ServerKeyExchange\\ + ServerHelloDone} (0,-2.7);
    \draw[panah] (0,-4.2) -- node[label]
        {ClientKeyExchange\\ + ChangeCipherSpec\\ + Finished} (5.6,-4.2);
    \draw[panah] (5.6,-5.4) -- node[label]
        {ChangeCipherSpec\\ + Finished} (0,-5.4);
    \draw[balik] (0,-6.6) -- node[label] {Data terenkripsi} (5.6,-6.6);

    % ================= TLS 1.3 (kanan) =================
    \begin{scope}[xshift=7.6cm]
    \node[font=\small\bfseries] at (2.8,0.95) {TLS 1.3: satu putaran};
    \node[kepala] (k13) at (0,0.35) {Klien};
    \node[kepala] (s13) at (5.6,0.35) {Server};
    \draw[hidup] (0,0) -- (0,-7.2);
    \draw[hidup] (5.6,0) -- (5.6,-7.2);

    \draw[panah] (0,-1.3) -- node[label] {ClientHello\\ + key share} (5.6,-1.3);
    \draw[panah] (5.6,-3.4) -- node[label]
        {ServerHello\\ + key share\\ \{Certificate,\\ CertificateVerify,\\ Finished\}} (0,-3.4);
    \draw[panah] (0,-5.4) -- node[label] {Finished} (5.6,-5.4);
    \draw[balik] (0,-6.6) -- node[label] {Data terenkripsi} (5.6,-6.6);
    \end{scope}

    \end{tikzpicture}
    \caption{Perbandingan jabat tangan TLS 1.2 dan TLS 1.3. TLS 1.2 menuntut
    dua putaran penuh sebelum data dapat dikirim, karena kunci sesi baru
    ditentukan setelah klien menerima sertifikat server. TLS 1.3 memangkasnya
    menjadi satu putaran: klien sudah mengirim bagian kunci publiknya pada
    pesan pertama, sehingga server dapat menghitung kunci sesi lebih awal dan
    seluruh pesan setelah \textit{ServerHello} sudah terenkripsi}
    \label{fig:jabat-tangan-tls}
\end{figure}
.
\begin{figure}[htbp]
    \centering
    \begin{tikzpicture}[
        panah/.style={-{Latex[length=2.2mm]}, thick},
        balik/.style={-{Latex[length=2.2mm]}, thick, dashed},
        label/.style={font=\scriptsize, align=center, fill=white, inner sep=1.5pt},
        kepala/.style={draw, rounded corners, fill=blue!6, font=\small,
                       minimum width=1.7cm, minimum height=0.6cm},
        hidup/.style={dashed, gray},
    ]

    % ================= TLS 1.2 (kiri) =================
    \node[font=\small\bfseries] at (2.8,0.95) {TLS 1.2: dua putaran};
    \node[kepala] (k12) at (0,0.35) {Klien};
    \node[kepala] (s12) at (5.6,0.35) {Server};
    \draw[hidup] (0,0) -- (0,-7.2);
    \draw[hidup] (5.6,0) -- (5.6,-7.2);

    \draw[panah] (0,-1.3) -- node[label] {ClientHello} (5.6,-1.3);
    \draw[panah] (5.6,-2.7) -- node[label]
        {ServerHello + Certificate\\ + ServerKeyExchange\\ + ServerHelloDone} (0,-2.7);
    \draw[panah] (0,-4.2) -- node[label]
        {ClientKeyExchange\\ + ChangeCipherSpec\\ + Finished} (5.6,-4.2);
    \draw[panah] (5.6,-5.4) -- node[label]
        {ChangeCipherSpec\\ + Finished} (0,-5.4);
    \draw[balik] (0,-6.6) -- node[label] {Data terenkripsi} (5.6,-6.6);

    % ================= TLS 1.3 (kanan) =================
    \begin{scope}[xshift=7.6cm]
    \node[font=\small\bfseries] at (2.8,0.95) {TLS 1.3: satu putaran};
    \node[kepala] (k13) at (0,0.35) {Klien};
    \node[kepala] (s13) at (5.6,0.35) {Server};
    \draw[hidup] (0,0) -- (0,-7.2);
    \draw[hidup] (5.6,0) -- (5.6,-7.2);

    \draw[panah] (0,-1.3) -- node[label] {ClientHello\\ + key share} (5.6,-1.3);
    \draw[panah] (5.6,-3.4) -- node[label]
        {ServerHello\\ + key share\\ \{Certificate,\\ CertificateVerify,\\ Finished\}} (0,-3.4);
    \draw[panah] (0,-5.4) -- node[label] {Finished} (5.6,-5.4);
    \draw[balik] (0,-6.6) -- node[label] {Data terenkripsi} (5.6,-6.6);
    \end{scope}

    \end{tikzpicture}
    \caption{Perbandingan jabat tangan TLS 1.2 dan TLS 1.3. TLS 1.2 menuntut
    dua putaran penuh sebelum data dapat dikirim, karena kunci sesi baru
    ditentukan setelah klien menerima sertifikat server. TLS 1.3 memangkasnya
    menjadi satu putaran: klien sudah mengirim bagian kunci publiknya pada
    pesan pertama, sehingga server dapat menghitung kunci sesi lebih awal dan
    seluruh pesan setelah \textit{ServerHello} sudah terenkripsi}
    \label{fig:jabat-tangan-tls}
\end{figure}
